% ReportKit acceptance test: presentation publication type / slides renderer
% / executive theme (Phase B). Exercises every reportkit-presentation.sty
% composition once, plus reuse of the shared semantic modules
% (reportkit-boxes, reportkit-code, reportkit-diagrams) already proven under
% the paged renderer -- this document exists to prove they also compile
% unmodified under Beamer, not to be a real deck. Compiled directly (like
% latex_templates/examples/equity-research/report.tex), not through the
% publication_pipeline -- see publication_pipeline/templates/slides-base.tex's
% header for why the pipeline path isn't wired up yet.
%
% Every composition is content-only (reportkit-presentation.sty's header
% explains why) and is wrapped here in an explicit \begin{frame}...\end{frame}
% -- [fragile] only where a frame's content actually needs it (codeblock).
\documentclass[theme=executive,publication-type=presentation]{reportkit-slides}

\setreportkitleftheader{ReportKit / Presentation Acceptance Test}
\setreportkitfooter{Presentation acceptance test}
\setreportkitversion{v0.1.0}
\title{Every presentation composition, once}
\subtitle{Phase B compile coverage, not a real deck}
\author{Test build}
\date{14 September 2026}

\begin{document}

\begin{frame}[plain]
\begin{titleslide}
\end{titleslide}
\end{frame}

\begin{frame}[plain]
\begin{sectiondivider}{Message and evidence}
\end{sectiondivider}
\end{frame}

\begin{frame}
\begin{messageslide}{Beamer renders every semantic box unmodified}
The callout, metric, code, and diagram modules were built against the renderer-hook
contract A3 introduced, not against the paged renderer specifically.
\end{messageslide}
\end{frame}

\begin{frame}
\begin{evidenceslide}{Three semantic modules compile under Beamer with zero changes}
\begin{principle}{Reused, not reimplemented}
\texttt{reportkit-boxes}, \texttt{reportkit-code}, and \texttt{reportkit-diagrams} are
required by \texttt{reportkit-slides.cls} exactly as \texttt{reportkit.cls} requires
them for the paged renderer.
\end{principle}
\end{evidenceslide}
\end{frame}

\begin{frame}[plain]
\begin{sectiondivider}{Visuals}
\end{sectiondivider}
\end{frame}

\begin{frame}
\begin{visualtext}[right]{A layered diagram}{RKLayer renders identically under
either renderer; only the placement hooks differ.}
\begin{diagram}[
  type=layer,
  caption={System layers, top to bottom.},
  source={Illustrative architecture.},
  description={Three stacked layers: presentation, application logic, and data storage.}
]
  \RKLayerSetup{7}{0.7}
  \RKLayer[rk accent]{0}{Presentation}{Renders slides for the reader.}
  \RKLayer{1}{Application logic}{Applies ReportKit macros.}
  \RKLayer{2}{Data storage}{Holds source content.}
\end{diagram}
\end{visualtext}
\end{frame}

\begin{frame}
\begin{fullvisual}[A metric card, unmodified under Beamer.][Illustrative.]
\begin{metric}{Compile coverage}{100\%}
Every composition in this fixture appears exactly once.
\end{metric}
\end{fullvisual}
\end{frame}

\begin{frame}[plain]
\begin{sectiondivider}{Structure}
\end{sectiondivider}
\end{frame}

\begin{frame}
\begin{comparison}
\comparisoncolumn{Paged}{Continuous page flow; \texttt{\string\RKReserveSpace} reserves
room before a break.}
\comparisoncolumn{Slides}{One composition per frame; \texttt{\string\RKReserveSpace} is a
no-op.}
\end{comparison}
\end{frame}

\begin{frame}
\begin{threepart}
\threepartcolumn{Registry}{Python is canonical; the generated \texttt{.def} is a
backstop.}
\threepartcolumn{Core}{Renderer-neutral mechanism and the style-token contract.}
\threepartcolumn{Renderer core}{Paged or slides; implements the same hook names.}
\end{threepart}
\end{frame}

\begin{frame}
\begin{herometric}{160mm x 90mm}{Declared canvas}
Beamer's own aspectratio=169 table already produces this frame size natively.
\end{herometric}
\end{frame}

\begin{frame}[fragile]
\begin{codeblock}[reportkit-slides.cls]{TeX}
\RequirePackage{reportkit-slides-core}
\RKAssertRendererHooks
\end{codeblock}
\begin{outputblock}
PASS -- renderer hooks loaded
\end{outputblock}
\end{frame}

\begin{frame}[plain]
\begin{appendixdivider}{Chart, table, and architecture aliases}
\end{appendixdivider}
\end{frame}

\begin{frame}
\begin{chartslide}[left]{Chart slide}{Same layout as visualtext, named for
discoverability.}
\begin{diagram}[
  type=matrix,
  caption={Illustrative placement.},
  description={A 2x2 matrix.}
]
  \RKMatrixSetup{6}{4}
  \RKMatrix{Effort}{Impact}
  \RKMatrixCell{nw}{Quick win}{Low effort, high impact}
  \RKMatrixCell{se}{Reconsider}{High effort, low impact}
\end{diagram}
\end{chartslide}
\end{frame}

\begin{frame}
\begin{tableslide}[right]{Table slide}{Same layout as visualtext, mirrored for a
table instead of a chart.}
\begin{metric}{Table stand-in}{4/4}
Compositions covered by this fixture so far.
\end{metric}
\end{tableslide}
\end{frame}

\begin{frame}
\begin{architectureslide}[An architecture diagram, full-bleed.][]
\begin{diagram}[
  type=layer,
  caption={Full-bleed variant.},
  description={Two stacked layers.}
]
  \RKLayerSetup{9}{0.9}
  \RKLayer[rk accent]{0}{Frontend}{Renders the interface.}
  \RKLayer{1}{Backend}{Serves data.}
\end{diagram}
\end{architectureslide}
\end{frame}

\begin{frame}[plain]
\begin{closingslide}{Thank you}
Questions welcome.
\end{closingslide}
\end{frame}

\end{document}
